\chapter{User experimentation forms}
\label{userExperForms}
\section{Consent form}
\label{sec:consent}
\begin{center}
{\LARGE \bf University of Cape Town} \vspace{10pt} \\ 
{\large \bf Department of Computer Science}
\end{center}
{\bf Consent form} \vspace{15pt} \\
Researcher: \line(90,0){150}

\begin{itemize}
\item I agree to participate in this experiment.
\item I agree to my responses being used for education and research.
\item I understand that my personal details will be used in aggregate form only, so that I will not be personally identifiable.
\item I understand that I am under no obligation to take part in this project.
\item I understand I have the right to withdraw from this experiment at any stage.
\item I have read this consent form and the information it contains and had the opportunity to ask questions about them.
\end{itemize}
\begin{tabbing}
Name of Participant: \indent \= \line(90,0){150}     \vspace{15pt} \\
Signature of Participant:  \> \line(90,0){150}  \vspace{15pt} \\
Date: \>  \line(90,0){150}        \vspace{10pt}       \\
\end{tabbing}



\newpage
\section{Introductory instruction form}
\label{introInstruct}

{\bf \large Gaming on the iPad introduction}
\parindent 0cm
\parskip 0.5cm

Thank you for agreeing to participate in this study. Please read this information carefully as it explains the experimental procedure.

The experiment will take approximately 45 minutes and will begin with a practice session where you will be asked to draw a number of gestures. Once this has been completed, you will play a game three times. Between each game session you will be asked to describe your game-play experience by completing a questionnaire. Please ensure that you fill out the questionnaire marked with the same number as instructed within the game. Further instructions on how to play the game will be given within the game. Please play each game until told to stop.

A demographic information form has been included with the questionnaires. Please fill this out before you begin the practice session.

If you have any further questions, please ask the available experimenter.
\parindent 1cm
\parskip 0.0cm

\newpage
\section{Demographic information form}
\label{demographic}
\begin{center}
{\large \bf Demographic information}
\end{center}

\begin{tabbing}
{\bf Age:} \indent \indent \= \line(90,0){150} \\ 
{\bf Gender:} \> \line(90,0){150} \\ 
{\bf Qualification towards which you are studying:} \line(90,0){150} \\
{\bf Year of study:} \> \line(90,0){150} \\
\end{tabbing}

{\bf Please rate your level of computer literacy:} \vspace{15pt}\\
\begin{center}
\begin{picture}(475,10)

  \put(0,20){\rule{450pt}{0.5pt}}
  \put(0,15){\rule{0.5pt}{10pt}}
  \put(75,15){\rule{0.5pt}{10pt}}
  \put(150,15){\rule{0.5pt}{10pt}}
  \put(225,15){\rule{0.5pt}{10pt}}
  \put(300,15){\rule{0.5pt}{10pt}}
  \put(375,15){\rule{0.5pt}{10pt}}
  \put(450,15){\rule{0.5pt}{10pt}}
  

  \put(-15,5){\footnotesize{Illiterate}}
  \put(425,5){\footnotesize{Highly literate}}
  

  \put(-3,30){1}
  \put(72,30){2}
  \put(147,30){3}
  \put(222,30){4}
  \put(297,30){5}
  \put(372,30){6}
  \put(447,30){7}
  
  
\end{picture}
\end{center}

\vspace{10pt}
{\bf Please rate how often you play games:} \vspace{15pt}\\
\begin{center}
\begin{picture}(475,10)

  \put(0,20){\rule{450pt}{0.5pt}}
  \put(0,15){\rule{0.5pt}{10pt}}
  \put(75,15){\rule{0.5pt}{10pt}}
  \put(150,15){\rule{0.5pt}{10pt}}
  \put(225,15){\rule{0.5pt}{10pt}}
  \put(300,15){\rule{0.5pt}{10pt}}
  \put(375,15){\rule{0.5pt}{10pt}}
  \put(450,15){\rule{0.5pt}{10pt}}
  

  \put(-10,5){\footnotesize{Never}}
  \put(440,5){\footnotesize{Daily}}
  

  \put(-3,30){1}
  \put(72,30){2}
  \put(147,30){3}
  \put(222,30){4}
  \put(297,30){5}
  \put(372,30){6}
  \put(447,30){7}
  
  
\end{picture}
\end{center}
\vspace{10pt}

{\bf Please rate your experience with touch-devices:} \vspace{15pt}\\
\begin{center}
\begin{picture}(475,10)

  \put(0,20){\rule{450pt}{0.5pt}}
  \put(0,15){\rule{0.5pt}{10pt}}
  \put(75,15){\rule{0.5pt}{10pt}}
  \put(150,15){\rule{0.5pt}{10pt}}
  \put(225,15){\rule{0.5pt}{10pt}}
  \put(300,15){\rule{0.5pt}{10pt}}
  \put(375,15){\rule{0.5pt}{10pt}}
  \put(450,15){\rule{0.5pt}{10pt}}
  

  \put(-20,5){\footnotesize{No experience}}
  \put(420,5){\footnotesize{Very experienced}}
  

  \put(-3,30){1}
  \put(72,30){2}
  \put(147,30){3}
  \put(222,30){4}
  \put(297,30){5}
  \put(372,30){6}
  \put(447,30){7}
  
  
\end{picture}
\end{center}

\newpage
\section{Flow State Scale}
\label{sec:flowScale}
Due to copyright limitations, only a select portion of the questionnaire is able to be reproduced. Sample questions are included below. \vspace{15pt}\\

{\bf I was challenged, but I believed my skills would allow me to meet the challenge} \vspace{15pt}\\
\begin{center}
\begin{picture}(325,10)

  \put(0,20){\rule{300pt}{0.5pt}}
  \put(0,15){\rule{0.5pt}{10pt}}
  \put(75,15){\rule{0.5pt}{10pt}}
  \put(150,15){\rule{0.5pt}{10pt}}
  \put(225,15){\rule{0.5pt}{10pt}}
  \put(300,15){\rule{0.5pt}{10pt}}
  
  

  \put(-30,5){\footnotesize{Strongly disagree}}
  \put(275,5){\footnotesize{Strongly agree}}
 

  \put(-3,30){1}
  \put(72,30){2}
  \put(147,30){3}
  \put(222,30){4}
  \put(297,30){5}
  
  
  
\end{picture}
\end{center}

\indent {\bf It was no effort to keep my mind on what was happening} \vspace{15pt}\\
\begin{center}
\begin{picture}(325,10)

  \put(0,20){\rule{300pt}{0.5pt}}
  \put(0,15){\rule{0.5pt}{10pt}}
  \put(75,15){\rule{0.5pt}{10pt}}
  \put(150,15){\rule{0.5pt}{10pt}}
  \put(225,15){\rule{0.5pt}{10pt}}
  \put(300,15){\rule{0.5pt}{10pt}}
  
  

  \put(-30,5){\footnotesize{Strongly disagree}}
  \put(275,5){\footnotesize{Strongly agree}}
 

  \put(-3,30){1}
  \put(72,30){2}
  \put(147,30){3}
  \put(222,30){4}
  \put(297,30){5}
  
  
  
\end{picture}
\end{center}
\vspace{10pt}

\indent {\bf I loved the feeling of what I was doing, and want to capture this feeling again} \vspace{15pt}\\
\begin{center}
\begin{picture}(325,10)

  \put(0,20){\rule{300pt}{0.5pt}}
  \put(0,15){\rule{0.5pt}{10pt}}
  \put(75,15){\rule{0.5pt}{10pt}}
  \put(150,15){\rule{0.5pt}{10pt}}
  \put(225,15){\rule{0.5pt}{10pt}}
  \put(300,15){\rule{0.5pt}{10pt}}
  
  

  \put(-30,5){\footnotesize{Strongly disagree}}
  \put(275,5){\footnotesize{Strongly agree}}
 

  \put(-3,30){1}
  \put(72,30){2}
  \put(147,30){3}
  \put(222,30){4}
  \put(297,30){5}
  
  
  
\end{picture}
\end{center}
\vspace{10pt}

\indent {\bf I had a feeling of total control over what I was doing} \vspace{15pt}\\
\begin{center}
\begin{picture}(325,10)

  \put(0,20){\rule{300pt}{0.5pt}}
  \put(0,15){\rule{0.5pt}{10pt}}
  \put(75,15){\rule{0.5pt}{10pt}}
  \put(150,15){\rule{0.5pt}{10pt}}
  \put(225,15){\rule{0.5pt}{10pt}}
  \put(300,15){\rule{0.5pt}{10pt}}
  
  

  \put(-30,5){\footnotesize{Strongly disagree}}
  \put(275,5){\footnotesize{Strongly agree}}
 

  \put(-3,30){1}
  \put(72,30){2}
  \put(147,30){3}
  \put(222,30){4}
  \put(297,30){5}
  
  
  
\end{picture}
\end{center}
\vspace{10pt}

\indent {\bf It felt like time went by quickly} \vspace{15pt}\\
\begin{center}
\begin{picture}(325,10)

  \put(0,20){\rule{300pt}{0.5pt}}
  \put(0,15){\rule{0.5pt}{10pt}}
  \put(75,15){\rule{0.5pt}{10pt}}
  \put(150,15){\rule{0.5pt}{10pt}}
  \put(225,15){\rule{0.5pt}{10pt}}
  \put(300,15){\rule{0.5pt}{10pt}}
  
  

  \put(-30,5){\footnotesize{Strongly disagree}}
  \put(275,5){\footnotesize{Strongly agree}}
 

  \put(-3,30){1}
  \put(72,30){2}
  \put(147,30){3}
  \put(222,30){4}
  \put(297,30){5}
  
  
  
\end{picture}
\end{center}
\vspace{10pt}


\newpage
\section{Post-experiment form}
\label{postEForm}

{\bf \large Post-experiment explanation}
\parindent 0cm
\parskip 0.5cm

Thank you for taking part in this experiment.

The research question being investigated was whether gestures make touch-based games more enjoyable and whether there are genre dependent differences in this.

In order to investigate this, three game version were designed. These included a gestural interface, direct manipulation interface and a version that contained a mix of the previous two. The questionnaire you answered assesses the level of flow present in each condition. Flow can be understood as the intrinsic enjoyment, or fun, one gets out of performing some activity. In this case, we're using flow to determine which game version was the most enjoyable to play.

In order to determine whether there are genre dependent differences in the appropriateness of gestures in gaming, three different games were developed. These included a strategy, puzzle and duelling/fighting game. Once we have all the results, comparisons will be done both across game version and genre.

Once again, thank you for participating in our experiment. If you have any questions, feel free to ask.

\parindent 1cm
\parskip 0.0cm
